\chapter{Proof of Stake \small{\textsf{DRAFT}}}\label{chapter.stake}

In \Cref{chapter.untrusty-world}, we established that money is a social construct:
Money is made valuable by a counter-party willing to accept it and hold it;
\emph{society at large} must be willing to accept it and hold it.
As a holder of money, we trust that the majority of society will act in this manner,
or, at the very least, those members of society who are economically active,
the \emph{economic majority}. It is only natural to assume that the majority of
coins in our system are in honest hands.

This is what gave rise to the assumption we made when we designed proof-of-work consensus,
namely that the majority of computational power is in honest hands. Since the majority of
money belongs to honest parties, and these parties are free to purchase computational power,
we deduce that the majority of computational power is honest. However, this conclusion is imperfect
for a couple of reasons. First, the adversary may have access to better hardware than the honest
parties, for example because they control a GPU production company. Second, even if the adversary
has access to the same hardware market, economics of scale apply; purchasing equipment
worth $10\times$ the cost gets her more than $10\times$ the computational power. Thus,
even if the adversary controls a minority of the money, if that money is concentrated,
she may be able to obtain a majority of the computational power.

Besides these concerns pertaining to the truthfulness of our assumptions,
proof-of-work has another major drawback: it is energy-intensive due to mining.
The more secure we want the network to be, the more energy we have to expend collectively
on mining. This energy expenditure is not per se ``wasted'' because it provides security.
Nevertheless, it would be nice to avoid it if we can find a way to obtain the same level
of security without it.

It seems that we have been using the honest computational majority assumption as an
imperfect \emph{proxy} for the underlying social assumption of money being in honest hands.
Ideally, we would like to directly use the belief that the majority of money is in honest hands
as our assumption. We call this the \emph{honest stake majority assumption}\index{Honest Stake Majority Assumption}.

\section{The Proof-of-Stake Inequality}

Recall that the proof-of-work inequality is $H(s \concat x \concat \ctr) \leq T$, where $s$ is the previous blockid,
$x$ is the commitment to the transactions, and $\ctr$ is the nonce. Brute forcing the nonce in this inequality
gives rise to rare events, which are necessary to ensure safety (e.g., avoid double spends). The more compute
one has, the more $\ctr$ values one can try, and the higher the probability of finding a ticket.
Our purpose now is to alter this inequality so that we still have rare events, but no brute forcing is
required, thereby bypassing the computational needs. Instead, we want the probability of finding a ticket
to depend on how much money one has.

When we constructed the execution layer, we invented two models that represent the money in the system:
the \emph{UTXO model} and the \emph{account model}. But here we have a chicken-and-egg problem:
To build the execution layer, we need the consensus layer, which requires the ticketing mechanism;
but to build a proof-of-stake ticketing mechanism, we need the execution layer that holds the information
of money ownership. The problem is even more pronounced: The very assumption that the stake majority
is honest rests upon a social belief pertaining to the coins produced at the execution layer.
We believe money is in honest hands; this belief gives a stake distribution; using this stake
distribution, we build a lottery; using this lottery, we create the money which we believed
to be in honest hands in the first place.
We illustrate this recursive dependency in \Cref{fig.pos-chicken-egg}.
We will put this problem aside for the moment and return to it later.
For now, we will assume that the population of money holders is fixed ahead of time, with
each member of that population holding exactly one unit of money. The distribution of money
across parties is called the \emph{stake distribution}\index{Stake Distribution}.
We also assume that this
population is known ahead of time, meaning that all parties have a list with the public keys
of all money holders at the beginning of the execution. This list of public keys is called a
\emph{Public Key Infrastructure}\index{Public Key Infrastructure} (PKI)\index{PKI}.

\begin{figure}[htbp!]
  \centering
  \includegraphics[width=0.7 \columnwidth,keepaspectratio]{14-stake/figures/pos-chicken-egg.pdf}
  \caption{The chicken-and-egg problem of proof-of-stake. The bottom layer is a social consensus
           assumption needed for money to be useful. Above it is our cryptographic assumption that
           the majority of coins are in honest hands, which is supported by the underlying
           social contract. Using this assumption, we construct the proof-of-stake consensus
           layer, on top of which an execution layer that creates money is built. This money
           is exactly the money that the bottom-layer social contract pertains to.}
  \label{fig.pos-chicken-egg}
\end{figure}

Let us take a first shot at altering the inequality. Because we have assumed our population
is static, the target $T$ will be a fixed constant. Let us remove the nonce $\ctr$ from the inequality,
since its sole purpose was to capture computational effort. We want each of our participants
to have an independent chance of finding a ticket, so we will add their public key $\pk$ to the
inequality.

\[
  H(s \concat x \concat \pk) \leq T
\]

If someone succeeds in finding a ticket, we call them the \emph{block proposer}\index{Block Proposer}.
Because we are now removing the computational burden to creating a block, we will speak of
\emph{proposing} blocks instead of \emph{mining} them.

Consider the beginning of the execution in which $s$ is set to the genesis block and $x$ is empty,
as no transactions have appeared in the mempool yet. Given this particular $s$ and $x$, an honest
party holding the public key $\pk$ has exactly one
shot at finding a ticket. We would like each corrupted party holding a public key to also get
exactly one shot. But nobody forces the adversary to fix their $x$! Given just one
corrupted public key $\pk$, the adversary can try brute-forcing multiple values for $x$
by locally trialing transactions of her own. This creates an unfair advantage for the adversary,
giving her a much higher probability of success compared to her stake. Somehow, the adversary
has utilized her computational power to beat the honest parties who are just using their stake.
When computational power is abused within stake systems, we call it a
\emph{grinding attack}\index{Grinding Attack}. Another way of mounting a grinding attack, when
the mempool is not empty, is to brute-force across valid permutations of the mempool $\overline{x}$.
Yet another way to grind is when the execution has already produced some past blocks.
The adversary can alter $s$ and choose to extend whichever fork enables her to satisfy the inequality.
To avoid grinding attacks, we must remove $s$ and $x$ from the inequality and obtain
\[
  H(\pk) \leq T\,.
\]

This inequality does not look very healthy. The first problem is that we have lost all information about the block
content. Where did the transactions and ancestry go? We will fix this issue by requiring
the proposer to accompany the ticket with a separate \emph{signed} block header.
At this time, we have separated the ticket from the block header.
The \emph{ticket} is a proof that the proposer is legitimate; for now, just their
public key. This ticket is verifiable by the rest of the parties (by redoing the hash and checking
the inequality) and allows them to know who the proposers are. Now, when the proposer issues a block,
they broadcast the tuple
\[
(\pk, (s \concat x \concat \pk \concat \sigma))\,,
\]
where $\pk$ is the ticket, $s \concat x \concat \pk \concat \sigma$ is the block header,
and $\sigma$ is the signature of the proposer. Upon receiving this tuple, the recipient checks
that:

\begin{enumerate}
  \item The ticket $\pk$ is valid by ensuring that $H(\pk) \leq T$.
  \item The signature $\sigma$ is valid by checking that $\Ver(\pk, (s \concat x \concat \pk), \sigma)$.
\end{enumerate}

Note that, naturally, the plaintext that the proposer signs is missing the signature.

The second problem is that our system is static. Instead,
we would like to create a dynamic system in which parties regularly alternate
between taking the proposer role as time passes. To do this, we will split time
into fixed-size intervals (in the order of a few seconds)
called \emph{slots}\index{Slot},
\glsxtrnewsymbol[description={slot}]{slot}{$\slot$}\glsadd{slot}
with the intent of
having approximately one block produced in each slot. We will number slots using
the natural numbers $\slot = 0, 1, 2, \ldots$. We want the slot to be longer than the
network delay so that a block honestly produced during one slot is seen by
all honest parties before the next slot begins. Contrary to rounds,
slot numbers will be used by our protocol, whereas rounds are an analysis
tool. But the two have a concrete relationship, which is that the slot duration
is at least $\Delta$ rounds.

Let us now add the slot number to the block header
\[
B = s \concat x \concat \slot \concat \pk \concat \sigma
\]
and to the proof-of-stake inequality
\[
H(\slot \concat \pk) \leq T\,.
\]

\section{Longest Chain Proof-of-Stake}

We have now a dynamic mechanism in which, for every slot, each party has a shot at being
the block proposer. The way our protocol works is exactly identical to the proof-of-work longest
chain protocol. Namely, there is a known genesis block $\mathcal{G}$.
The PKI is encoded within the genesis block.
Each honest party maintains a local chain and mempool. Upon receiving
a new chain from the network, the party validates it and adopts it if it is longer than the local chain.
Validation involves:

\begin{enumerate}
  \item Check that the first block is the genesis block.
  \item Check the ancestry of each block.
  \item Check that the $\pk$ signing the block is in the PKI.
  \item Check the proof-of-stake inequality for the given slot and $\pk$.
  \item Check that the tip has a slot not in the future.
  \item Check that the block slots are strictly increasing.
\end{enumerate}

To propose a new block,
at the beginning of each slot, the party checks the proof-of-stake inequality to determine if they
are a block proposer. If so, they create a new block using their local mempool and point it to the
tip of their current local chain. They sign this block and broadcast it. Identically to the
proof-of-work longest chain protocol, a block is considered
\emph{stable} when it is buried under $k$ blocks, and a transaction becomes \emph{confirmed} when it is part
of the stable chain.

We observe that, similar to proof-of-work, within a single slot there can be zero, one, or more
proposers. We prefer the case where there is exactly one honest proposer per slot which, following
our terminology from \Cref{chapter.chain}, we name a \emph{convergence opportunity}.
When there are multiple
honest proposers per slot, we can have fan-out situations. These will eventually be resolved for the
same reason they are resolved in proof-of-work: At some point, there will be a convergence
opportunity which will not be matched by an adversarial block. Naturally, this requires that
the majority of stake is honest.

% TODO: figures about past and future
During validation, we demand that slots in blocks are not in the future. This requires the
honest party to keep a local clock and compare this local clock against the slot numbers in the
blocks they receive. If a block has a slot number in the future, it is rejected. If this check
was missing, the adversary could test the proof-of-stake inequality for a massive number
of future slots \emph{in advance} and leverage
all these tickets to create a chain \emph{today}, whereas the honest parties only test the
current slot.
This highlights a fundamental difference between
proof-of-work and proof-of-stake: In proof-of-work, the parties cannot predict future successes
until the time of the success actually occurs, exactly because the expenditure of computational
power takes time. On the other hand, in proof-of-stake, a party can peek into the future and
know in which slots she will be a proposer. This is the reason why we need the check that
slots of valid blocks are only in the past.

Dually, we also demand that the slots of the blocks in a chain are increasing.
This prevents the adversary from using old blocks in the future. If the check was missing, the following
concrete attack would be possible: The adversary sleeps for a long period of time and allows the honest parties
to produce a chain $\chain$. At some point in time, the adversary wakes up and uses all of her past
tickets on top of $\chain[{:}{-k}]$, breaking safety. In proof-of-work, we avoided such issues
by including the previd $s$ in the ticketing mechanism, ensuring the proof-of-work is invalidated whenever the
$s$ contained within is retroactively changed.

Contrary to proof-of-work, in proof-of-stake,
a successful ticket does not bind the adversary to a particular block.
Instead, the adversary may use the same ticket to produce multiple blocks
for the same slot. This is done by simply sending the same ticket with
multiple signatures over different block headers, potentially to different
honest recipients, splitting the honest population between different chain tips.
We call this adversarial behavior an \emph{equivocation}\index{Equivocation}.
The reason equivocations are possible is because in proof-of-stake, there
is no additional cost associated with signing more blocks per ticket.
In a sense, there is \emph{nothing at stake}\index{Nothing at Stake}.
On the other hand, in the proof-of-work mechanism,
each new block must be accompanied by a novel
proof of computational work. Producing that proof of work has a cost
that cannot be recovered or reused in the future.

In \Cref{fig.pos-example}, we show an example of a proof-of-stake execution.
Slot numbers are shown as integers at the top.
Each slot is marked as honest $H$ if at least one honest party is a proposer for that slot;
as adversarial $\mathcal{A}$ if at least one corrupted party is a proposer for that slot;
or as $\bot$ if there is no proposer for that slot. Note that multiple parties can be proposers
for the same slot, hence the case $H / \mathcal{A}$ is possible (slot $8$).
Slot $0$ is marked as $H$ because the genesis block $\mathcal{G}$ is honest by convention.
Blocks are shown as blue rectangles and numbered based on their slot number.
Tickets are shown as circles surrounding the blocks that they accompany.
Honest tickets are shown in green dotted circles, whereas adversarial tickets are
shown in red solid circles. In the honest case, each ticket is associated with exactly one block.
The adversary may produce multiple blocks for the same slot by equivocating (slot $5$),
or may choose not to produce any block for a given ticket (slot $9$).
Note that in slot $7$ there are two honest proposers, each producing a different block
($7$a and $7$b). Each of these blocks extends a different parent chain of the same
length. At the conclusion of this execution, there is a unique longest chain known to all
honest parties. This chain has height $8$ and ends in block $10$ honestly produced in slot
$10$.

\begin{figure}[htbp!]
  \centering
  \includegraphics[width=1 \columnwidth,keepaspectratio]{14-stake/figures/pos-example.pdf}
  \caption{An example of a proof-of-stake execution. }
  \label{fig.pos-example}
\end{figure}

\section{Aggregation Independence}

Let us now remove the restriction that each party holds exactly one unit of money,
and allow different parties to hold different monetary amounts. The stake distribution
will now be a map from public keys.

\begin{definition}[Stake Distribution]
  A \emph{stake distribution} is a function $\SD\colon \PKI \longrightarrow [0, 1]$
  such that $\sum_{\pk \in \PKI} \SD(\pk) = 1$.
\end{definition}

To simplify later developments, we have normalized the stake distribution function
to the range $[0, 1]$, but, naturally, a real system must continue to use integer amounts
to avoid floating-point errors, a detail we are now abstracting away. We call this the
\emph{relative stake}\index{Relative Stake} of a party.

The honest majority
assumption needs to be altered to weight keys by the amount of money they hold:

\begin{definition}[Honest Stake Majority Assumption]
  For a given \emph{stake distribution} $\SD\colon \PKI \longrightarrow [0, 1]$
  of some public key infrastructure $\PKI$ among which $\Hon \subseteq \PKI$ are honest, we say
  that the \emph{honest stake majority assumption} holds if
  \[
    \sum_{\pk \in \Hon} \SD[\pk] > \sum_{\pk \in \PKI \setminus \Hon} \SD[\pk]\,.
  \]
\end{definition}

Our protocol should reflect this \emph{weighted} honest majority assumption by assigning
a higher probability of success to parties with more stake. Simply put, parties with
a higher stake should become proposers more often. We control this probability by altering
the target, so we will now make the target potentially different for each party, and
dependent on their stake. We define a function $\phi$ that, given a stake, returns
the target.

% TODO: define notation:
%  a) ranges [i:j], [i], [-j], [:-j], ...
%  b) the set [n]
%  c) empty string: $\epsilon$
%  d) Chernoff error parameter: $\varepsilon$
%  e) \delta \delta \delta

\begin{definition}[Target Function]
  A \emph{target function} is an increasing function $\phi\colon [0, 1] \longrightarrow [0, 1]$ that maps
  a normalized stake value to a normalized target, with $\phi(0) = 0$.
\end{definition}

The inequality then uses a different target per party:

\[
  \frac{H(\sl \concat \pk)}{2^\kappa} < \phi(\SD(\pk))
\]

Note that we have normalized the output of the hash value to the range $[0, 1]$ as well.
You can think of $\phi$ as the desired probability that a party with stake $\SD(\pk)$
obtains a block for a given slot.

The first function $\phi$ that comes to mind is the identity function $\phi(x) = x$.
Using this function, the probability of finding a block is exactly equal to the propotion
of the total stake that one owns. For this $\phi$, consider the case $n = 100$ and $t = 50$
in which the adversary controls exactly half of the parties, and suppose each party controls
the same amount of stake. Let us walk through the probability a given slot is honest.
This happens when at least one honest party succeeds. Namely, it is \emph{not}
the case that \emph{all} honest parties \emph{fail} to have a successful query in that slot.
\[
  \Pr[\text{the slot is honest}] = 1 - (1 - \frac{1}{50})^{50} \simeq 0.395\,.
\]
The probability is the same for an adversary that plays honestly.

The expected number of honest blocks
produced within the given slot is
\[
  \E[\text{blocks produced}] = 50 \times \frac{1}{100} = \frac{1}{2}\,,
\]
and the same is the case for the adversary playing honestly.

However, there is an adversarial
strategy that can improve her probability of getting an adversarial slot: The adversary
\emph{aggregates} the stake of her $50$ corrupted parties into a single public key, but
otherwise plays honestly. In
that case, her single party has a relative stake of $\frac{1}{2}$.
The expected number of adversarial blocks produced within the given slot remains $\frac{1}{2}$,
but the probability of getting an
adversarial slot is $\frac{1}{2} > 0.395$. On the other hand, honest parties do not have the ability
to aggregate their stake, and so are at a disadvantage.
This discrepancy of probabilities stems from the fact that, if multiple honest parties succeed in
finding a block in the same slot, this does not help the protocol make more progress than a single party
succeeding.

We'd like to design a target function $\phi$ which is \emph{aggregation independent}\index{Aggregation Independence}.
Namely the probability of succeeding in a given slot must be the same for the two following cases:
(a) for a \emph{group} of keys controlling a \emph{total} relative stake $\alpha$, and
(b) for a \emph{single} key aggregately controlling the same relative stake $\alpha$.

% TODO: broken example... then go to aggregation independence

\input{14-stake/algorithms/alg.pos-round}
\input{14-stake/algorithms/alg.pos-valid}

{\color{red}
\begin{itemize}
\item Combining honest and adversarial randomness to obtain unbiased randomness
\item Pseudorandomness: Getting more randomness from little randomness
\item Epoch randomness
\item Slot leaders
\item k+1 blocks in any continuous 2k slots
\item Verifiable random functions
\item The proof-of-stake inequality
\item The Ouroboros Praos protocol
\item Aggregation independence
\item Additional reading: egalitarianism
\end{itemize}
}

\section*{Problems}
\begin{problems}
  \item Consider the longest chain proof-of-stake protocol we studied in class. We will
    look at an instantiation where $\Delta = 1$ second, $n = 100$, $\kappa = 128$,
    $T = 2^{118}$. For answering the questions below, fix one constant $k$ of Your
    choice for the confirmation rule.

    \begin{enumerate}
        \item First suppose all $100$ nodes are honest.
            \begin{enumerate}
                \item Compute the expected growth rate of the longest
                    chain, in blocks per second.
                \item Compute the expected growth rate of the
                    $k$-confirmed chain, in blocks per second. (The $k$-confirmed
                    chain is the portion $C[{:}{-}k]$.)
                \item If we double $T$, does the expected growth rate of
                    the longest chain double, more than double, or less than double?
                    What about the growth rate of the $k$-confirmed chain?
            \end{enumerate}
        \item Now suppose $20$ of the $100$ nodes are adversary, the other $80$ are
            honest.
            \begin{enumerate}
                \item Compute a tight lower bound on the expected growth
                    rate of the longest chain. (A ``lower bound'' means that it is a
                    lower bound irrespective of the adversary's attack strategy;
                    ``tight'' means that the lower bound is attainable for some
                    adversary's attack strategy.)
                \item Compute a tight lower bound on the expected growth
                    rate of the $k$-confirmed chain.
                \item Is the protocol safe? Is the protocol live?
            \end{enumerate}
        \item Now suppose $20$ nodes are still adversary, but instead of having the
            full $80$ honest nodes online, $30$ of them decide not to participate in
            the protocol and went on vacation to Ikaria. However, the protocol
            designer does not know this and the protocol parameters are not
            adjusted.
            \begin{enumerate}
                \item Compute a tight lower bound on the expected growth
                    rate of the longest chain.
                \item Compute a tight lower bound on the growth rate of
                    the $k$-confirmed chain.
                \item Is the protocol safe? Is the protocol live?
            \end{enumerate}
        \item Now suppose a further $35$ honest nodes went on vacation.
            \begin{enumerate}
                \item Compute a tight lower bound on the expected growth
                    rate of the longest chain.
                \item Compute a tight lower bound on the expected growth
                    rate of the $k$-confirmed chain.
                \item Is the protocol safe? Is the protocol live?
            \end{enumerate}
        \item A protocol is said to be {\em available} if it is safe and
            live whenever the number of honest nodes online is greater than the
            number of adversary nodes. Is the longest chain PoS protocol available?
    \end{enumerate}
\end{problems}

\section*{Further Reading}

The first provably secure longest-chain proof-of-stake protocol is Ouroboros~\cite{ouroboros}.
This is used in practical systems such as Cardano and Polkadot. The first
construction that uses the proof-of-stake inequality in the way we defined in this chapter
appears in the Sleepy Model of Consensus~\cite{sleepy-model}. The proper definitions of
VRFs for use in the proof-of-stake inequality appear in Ouroboros Praos~\cite{praos},
where the notion of aggregation independence is also introduced.

Before these provably secure constructions, the idea of a proof-of-stake mechanism was
discussed in the Bitcoin community since 2011, with multiple unproven attempts to implement it,
with PPCoin being a notable example. The multitude of those attempts is illustrative of the
non-triviality of designing a secure longest-chain proof-of-stake protocol.